\section*{Introduction}

Taxonomic revisions turn observations into species descriptions, diagnoses, identification keys, and nomenclatural acts. Yet their preparation is still largely document-based: observations are rewritten as prose, edited in a manuscript, and reformatted for publication, typically in Microsoft Word or Google Docs. Repeatedly moving information between data and text takes time, makes it harder to keep them consistent, and can introduce errors. It also leaves much of the information in taxonomic publications difficult to find and reuse \citep{smith2013beyond, miller2012cybertaxonomic, penev2017strategies}.

This problem is especially clear in species descriptions. Natural language is easy for taxonomists to read, but computers cannot reliably identify the anatomical characters and their relationships described in the text. The same structure may also be named differently by different authors, while the meaning of a term may depend on the taxonomic group or context. These observations must therefore be extracted and recoded before they can be searched, compared across taxa, or combined with other data.

Semantic descriptions based on ontologies address this limitation by representing observations with standardized terms and explicit relationships \citep{balhoff2013semantic, deans2015phenotypes}. In simple terms, an ontology provides agreed names and definitions for biological structures and specifies how those structures are related. For example, it can distinguish an anatomical structure from one of its qualities and record that the structure is part of a particular organism. Each term has a stable identifier, allowing descriptions that use the same concept to be linked even when their wording differs. Phenoscript provides a formal language for writing such descriptions and can generate both machine-readable data and human-readable text for a conventional taxonomic treatment \citep{montanaro2024beyond, montanaro2024computable}. Taxonomists can therefore read the resulting description in a familiar form, while its underlying content remains available for computational searches and comparisons.

Generating readable text, however, does not solve the problem if it must be copied into the manuscript again after every change. We therefore connect Phenoscript, GitHub, and Overleaf in a single workflow to streamline this process. Phenoscript is the source of the species descriptions, Git records changes to the descriptions and related files, and Overleaf compiles the \LaTeX{} manuscript containing the generated text. The descriptions are written once in Phenoscript and rendered directly in the manuscript. When a description is corrected, its output can be regenerated and synchronized through GitHub instead of being copied, reformatted, and checked again by hand. This reduces the risk that the published treatment differs from its semantic source.

The same connection also makes the wider manuscript easier to manage. GitHub keeps the descriptions, manuscript, bibliography, figures, and processing files together and preserves a history of their changes. Collaborators can see what was changed, compare versions, and recover earlier content. Overleaf provides a shared writing environment and automatically compiles these files into the formatted paper. This approach keeps the manuscript closer to its source data and makes collaboration more transparent and reproducible \citep{ram2013git, braga2023github, himmelstein2020open}. Because \LaTeX{} separates content from formatting, the same source can also be adapted to journal templates with less manual reformatting \citep{bar2021reproducible, krewinkel2017formatting}.

Here, we use this workflow in a taxonomic review of the Afrotropical dung beetle genus \textit{Amietina} Cambefort, 1981 (Coleoptera: Scarabaeinae: Onthophagini). The genus currently contains four described species from forests in West, Central, and eastern Africa \citep{cambefort1981amietina, branco1989deux}. Material collected in the Eastern Arc Mountains of Tanzania revealed four additional species, broadening the morphological limits of the genus. The distributions of these species also add to the limited knowledge of dung beetle endemism across the Eastern Arc mountain blocks \citep{burgess2007biological, montanaro2024microallopatric}.

We describe \textit{A. bondarevi} \textbf{n. sp.}, \textit{A. grebennikovi} \textbf{n. sp.}, \textit{A. mazumbaiensis} \textbf{n. sp.}, and \textit{A. praedicta} \textbf{n. sp.}; describe the previously unknown female of \textit{A. nyassalandica}; provide new distributional records; and present an updated identification key. Alongside these taxonomic contributions, we show how Phenoscript, GitHub, and Overleaf can connect semantic phenotype data directly to a publication-ready manuscript.
